\chapter{Conclusiones}
\label{sec:conclusiones}

Esta tesis partió de una hipótesis concreta: que la propia señal de color de una nube de puntos
contiene información suficiente para construir transformadas gráficas más compactas que las que
se obtienen de la geometría sola, y que esa información puede aprovecharse sin que el costo de
comunicarla al decodificador anule la ganancia. El objetivo general fue, en consecuencia,
desarrollar un algoritmo determinístico de compresión de atributos, dependiente de la señal y
superior a su versión puramente volumétrica. Ese objetivo se cumple en su sentido estricto.
PCADC -- un libro de códigos direccional entrenado fuera de línea mediante k-means esférico
sobre $\mathbb{S}^2$, auto-lazos calibrados por optimización tasa-distorsión, selección de modo
por bloque padre mediante \textit{gating} y trellis de Viterbi, y señalización real del índice
de modo -- reduce el BD-Rate en $-2.18\%$ promedio sobre las 14 secuencias evaluadas (rango
$-1.39\%$ a $-3.91\%$ según la secuencia, Sección~\ref{sec:e10_bd_summary}) frente a la misma
transformada gráfica sin auto-lazos, sin ningún punto de operación en que sea peor que ella, y
con un costo de señalización que nunca supera el $1.09\%$ del flujo de bits. Es un algoritmo
determinístico, exactamente invertible, cuyo decodificador reconstruye grafo y transformada sin
acceso a la señal de color. Toda la evidencia es una comparación contra la propia línea base
estructural, no contra un códec de referencia externo; la ganancia demostrada es la del
mecanismo de adaptación por señal, y esa es la afirmación que este trabajo sostiene.

Los análisis del Capítulo de Resultados permiten caracterizar esa ganancia más allá del
promedio. En primer lugar, depende de la densidad de puntos y no de la etiqueta de profundidad:
las secuencias de profundidad 9 ganan sistemáticamente más ($-2.98\%$) que las de profundidad
10 ($-1.74\%$), la asimetría persiste al comparar a razón padre:hijo equivalente
(Sección~\ref{sec:e10_scale_ratio}), y el tamaño de la nube predice el desempeño mejor que la
profundidad nominal (Sección~\ref{sec:e10_driver_corr}). En segundo lugar, la ganancia crece
con la capacidad $K$ del libro de códigos en promedio y sin saturar hasta $K=24$, pero no de
forma uniforme: cerca de la mitad de las secuencias alcanza su óptimo en $K=24$ y no en $K=30$,
y una capacidad mayor amplifica la dispersión entre \textit{frames}
(Sección~\ref{sec:e10_k_sweep}). En tercer lugar, el overhead de señalización, además de
acotado, es informativo: es el correlato más fuerte de la ganancia ($r=-0.88$), de modo que las
secuencias que más gastan en señalizar son las que más ganan, y su comportamiento frente al
paso de cuantización se explica directamente por la barrera de \textit{gating}
(Sección~\ref{sec:mode_selection}). Por último, el entrenamiento se comporta como se esperaba:
el k-means esférico converge en todos los casos, las direcciones dominantes son estables entre
capacidades, y la calibración de auto-lazos extrae ganancia adicional con un número fijo y
pequeño de evaluaciones (Sección~\ref{sec:e10_codebook_dirs}). La inspección cualitativa es
coherente con lo anterior: no hay artefactos visibles y el error se redistribuye sobre la
superficie según la orientación del gradiente local, más que reducirse uniformemente
(Sección~\ref{sec:e10_qualitative}).

La observación más importante del trabajo concierne al compromiso entre adaptación y
señalización, y explica la distancia entre el techo teórico y el sistema real. Los oráculos de
la Sección~\ref{sec:oracle_upper_bounds} muestran que el mecanismo de auto-lazos es más
poderoso mientras más fina es la partición: hasta $-17.31\%$ en bloques de lado 4, con
conocimiento perfecto de la señal y sin señalizar nada. El sistema completo obtiene su mejor
desempeño en el régimen opuesto, con bloques hijo de lado 16, y en bloques de lado 4 apenas
supera la línea base. La razón no es que la adaptación deje de funcionar en bloques pequeños,
sino que una única decisión por bloque padre -- necesaria para acotar la señalización -- se
hereda a hijos cada vez más lejanos de la escala en que se tomó. El método captura hoy una
fracción del techo porque paga la señalización con resolución de decisión, y esa distancia
entre el nivel en que se decide y el nivel en que se transforma es su principal limitante y,
al mismo tiempo, el lugar donde queda la ganancia por encontrar. Dentro de ese compromiso,
varias decisiones de diseño quedan respaldadas por los datos: el \textit{gating} previo por
varianza reduce la señalización sin afectar la ganancia entre secuencias; la continuidad
espacial de los modos, explotada por el trellis y por el contexto del codificador de índices,
es una redundancia real; y la asimetría con que se calibran los auto-lazos -- una malla gruesa
para la fracción de nodos, un refinamiento fino solo para el peso -- coincide con la
distribución bimodal de los parámetros óptimos observada en los oráculos.

En lo metodológico, tres decisiones resultaron determinantes. La primera fue desacoplar el
entrenamiento: la aproximación inicial, una optimización evolutiva conjunta de dirección y
auto-lazos en cinco dimensiones, se reemplazó por un esquema en dos etapas
(Sección~\ref{sec:training_stage}) que es más simple, mucho más barato y, sobre todo,
evaluable a la escala de 14 secuencias, algo que la optimización conjunta nunca alcanzó; la
complejidad adicional no era necesaria para este problema. La segunda fue medir siempre con
codificación entrópica real, incluyendo el overhead, y sobre la curva tasa-distorsión completa
en vez de en un punto: solo así se distingue una ganancia de BD-Rate consistente de una
diferencia de PSNR que, a paso de cuantización fijo, es una fracción de decibel. La tercera fue
acotar antes de construir: los oráculos establecieron que existía ganancia que perseguir y
revelaron que una aproximación plana de la luminancia guía mejor al grafo que la luminancia
misma, lo que orientó todo el diseño del libro de códigos direccional.

Las limitaciones que corresponde reconocer son las que se derivan de las decisiones de diseño.
La decisión de modo se toma una vez por bloque padre y se hereda a todos sus hijos, y el tamaño
del padre queda fijado por la profundidad de la secuencia; esa elección acota la señalización
al costo ya descrito. El libro de códigos se entrena sobre un único \textit{frame} y por
secuencia, y su capacidad es discreta ($K\leq30$), con la fracción de auto-lazos restringida a
tres valores; se gana simplicidad y reproducibilidad a cambio de una aproximación de la
dirección de gradiente que deja ganancia sobre la mesa. Solo se codifica el canal de luminancia,
una práctica estándar que aquí deja fuera de la evaluación a la crominancia. Y la calibración
de los auto-lazos es costosa en secuencias densas -- entre 5.5 y 50 horas por secuencia en
profundidad 10, frente a menos de 30 minutos en profundidad 9
(Sección~\ref{sec:e10_multiframe_compute}) --, un costo que se paga una sola vez pero que
condiciona la escala a la que el método puede evaluarse.

Esas mismas limitaciones señalan dónde el método aún tiene que encontrar su mejor desempeño. La
línea más prometedora es una señalización multirresolución, en la que solo los hijos que lo
justifiquen por costo reciban un modo propio, acercando el nivel de decisión al nivel de
transformada con el margen de overhead disponible; le siguen aislar cuánto de la ganancia
proviene de la diversidad direccional frente al auto-lazo en sí, situar el resultado en una
escala externa mediante un códec de referencia, y verificar la generalización del libro de
códigos sobre secuencias completas y entre secuencias del mismo sujeto. La hipótesis inicial se
sostiene: la señal de color puede guiar la estructura del grafo con un costo de señalización
despreciable y producir una ganancia consistente sobre la transformada puramente estructural.
La contribución de esta tesis no está en la magnitud de esa ganancia, que es modesta frente al
techo que los oráculos demuestran alcanzable, sino en haber construido y medido un mecanismo
interpretable que convierte información de la señal en compresión, y en haber identificado con
precisión dónde se pierde el resto.
